%%% The lexeme hierarchy, shared by lec-08a-lexicon and lec-08b-lexrules.
%%% The lexeme hierarchy is shown repeatedly, with an arrow pointing
%%% at the type under discussion.  #1: comma-separated names of nodes
%%% to colour red; #2: TikZ code drawn after the tree (the arrow).
\makeatletter
\newcommand{\lxhl}[2]{\ifcsname lexhl@#1\endcsname\textcolor{red}{#2}\else #2\fi}
\newcommand{\lexhier}[2]{%
\begingroup
\@for\lx@n:=#1\do{\expandafter\def\csname lexhl@\lx@n\endcsname{}}%
\begin{center}
\scalebox{0.7}{\begin{forest}
cat, for tree={font=\itshape, s sep=3mm}
[{\lxhl{synsem}{synsem}\\\textnormal{\avm{[ \textsc{syn, sem} ]}}}, align=center, name=synsem, s sep=4mm
  [{\lxhl{lexeme}{lexeme}\\\textnormal{\avm{[ arg-st ]}}}, align=center, name=lexeme
    [{\lxhl{infl}{infl-lxm}}, name=infl
      [{\lxhl{verb}{verb-lxm}}, l*=3, name=verb
        [{\lxhl{siv}{siv-lxm}}, name=siv]
        [{\lxhl{piv}{piv-lxm}}, name=piv]
        [{\lxhl{tv}{tv-lxm}}, name=tv
          [{\lxhl{stv}{stv-lxm}}, name=stv]
          [{\lxhl{dtv}{dtv-lxm}}, name=dtv]
          [{\lxhl{ptv}{ptv-lxm}}, name=ptv]]]
      [{\lxhl{cn}{cn-lxm}}, l*=3, name=cn
        [{\lxhl{cntn}{cntn-lxm}}, name=cntn]
        [{\lxhl{massn}{massn-lxm}}, name=massn]]]
    [{\lxhl{const}{const-lxm}}, name=const
      [{\lxhl{adj}{adj-lxm}}, l*=2, name=adj]
      [{\lxhl{conj}{conj-lxm}}, l*=2, name=conj]
      [{\lxhl{det}{det-lxm}}, l*=2, name=det]
      [{\lxhl{predp}{predp-lxm}}, l*=2, name=predp]
      [{\lxhl{argmkp}{argmkp-lxm}}, l*=2, name=argmkp]
      [{\lxhl{pn}{pn-lxm}}, l*=1.4, name=pn]
      [{\lxhl{pron}{pron-lxm}}, l*=1.4, name=pron]]]
  [{\lxhl{expression}{expression}}, name=expression, l*=0.8
    [{\lxhl{word}{word}\\\textnormal{\avm{[ arg-st ]}}}, align=center, l*=0.6, name=word]
    [{\lxhl{phrase}{phrase}}, l*=0.6, name=phrase]]]
#2
\end{forest}}
\end{center}
\endgroup}
\makeatother
%%% the orange arrow
\tikzset{hlarrow/.style={overlay, ->, >=stealth, orange, line width=2mm}}
